\notechapter{N-20230311}{Riemann Integration and Lebesgue Integration: From Partitions of the Domain to Partitions of Values}

\section{Why two theories of integration exist}

This note compares Riemann integration and Lebesgue integration.  A compressed summary such as “Riemann cuts the $x$-axis; Lebesgue cuts the $y$-axis” is memorable, but it is not enough.  The real issue is what kinds of limiting processes the integral is supposed to tolerate.

Riemann integration works beautifully for continuous functions and many functions with mild discontinuities.  But once functions oscillate too much, especially when rational and irrational points behave differently, the Riemann approach becomes rigid.  Lebesgue integration was designed to handle limits of functions more naturally, especially in analysis, probability, and Fourier theory.

\section{Riemann sums}

Let $f:[a,b]\to\mathbb R$ be bounded.  A partition is a finite sequence
\[
  P:a=x_0<x_1<\cdots<x_n=b.
\]
Choose a sample point $\xi_i\in[x_{i-1},x_i]$ in each subinterval.  The Riemann sum is
\[
  \sum_{i=1}^n f(\xi_i)\Delta x_i,
  \qquad
  \Delta x_i=x_i-x_{i-1}.
\]
If all such sums approach the same limit as the mesh
\[
  \|P\|=\max_i \Delta x_i
\]
tends to zero, that limit is the Riemann integral.

The geometric picture is rectangles under a graph.  But the definition is not simply about drawing rectangles; it says that the total contribution should stabilize no matter how the sample points are chosen, once the partition is sufficiently fine.

\section{Darboux upper and lower sums}

A cleaner equivalent formulation is Darboux integration.  On each interval $[x_{i-1},x_i]$, define
\[
  M_i=\sup_{x\in[x_{i-1},x_i]} f(x),
  \qquad
  m_i=\inf_{x\in[x_{i-1},x_i]} f(x).
\]
The upper and lower sums are
\[
  U(P,f)=\sum_i M_i\Delta x_i,
  \qquad
  L(P,f)=\sum_i m_i\Delta x_i.
\]
Then
\[
  f\text{ is Riemann integrable}
  \quad\Longleftrightarrow\quad
  \inf_P U(P,f)=\sup_P L(P,f).
\]
The common value is the integral.

This formulation explains why discontinuities matter.  If a function oscillates wildly on every interval, then $M_i-m_i$ remains large no matter how small the interval is, and the upper and lower sums cannot meet.

\section{A simple integrable discontinuity}

Consider
\[
  f(x)=\begin{cases}
  1,&x=c,\\
  0,&x\ne c
  \end{cases}
\]
on $[a,b]$.  This function is discontinuous at one point, but it is Riemann integrable with integral zero.  Indeed, the lower sums are all zero.  For any partition, only one subinterval can contain $c$, and the contribution to the upper sum from that subinterval can be made arbitrarily small.

Thus Riemann integration tolerates finite discontinuities.  More generally, a bounded function is Riemann integrable iff its set of discontinuities has Lebesgue measure zero.  This theorem already hints that measure theory is hiding behind Riemann integration.

\section{Dirichlet's function}

The Dirichlet function is
\[
  1_{\mathbb Q}(x)=\begin{cases}
  1,&x\in\mathbb Q,\\
  0,&x\notin\mathbb Q.
  \end{cases}
\]
On every nonempty interval there are rational and irrational numbers.  Therefore on every subinterval,
\[
  M_i=1,
  \qquad
  m_i=0.
\]
Hence
\[
  U(P,f)=b-a,
  \qquad
  L(P,f)=0
\]
for every partition.  The upper and lower integrals never meet, so the function is not Riemann integrable.

Lebesgue integration handles this example immediately: since $\mathbb Q$ is countable and has measure zero,
\[
  \int_{[a,b]}1_{\mathbb Q}\,dx=0.
\]
This one example explains much of the motivation for Lebesgue's theory.

\section{The Lebesgue idea}

Riemann integration partitions the domain into small intervals and asks how the function behaves on each interval.  Lebesgue integration partitions according to values.  For a nonnegative simple function
\[
  s=\sum_{j=1}^m a_j1_{E_j},
\]
where the $E_j$ are measurable and disjoint, define
\[
  \int s\,d\mu=\sum_{j=1}^m a_j\mu(E_j).
\]
For a nonnegative measurable function $f$, define
\[
  \int f\,d\mu=\sup\left\{\int s\,d\mu:s\le f,\ s\text{ simple}\right\}.
\]
For a general integrable function, write
\[
  f=f^+-f^-.
\]
If both $\int f^+$ and $\int f^-$ are finite, define
\[
  \int f=\int f^+-\int f^-.
\]

The construction is less pictorial than Riemann sums, but more robust.  It counts how much space is occupied by each range of values.

\section{Characteristic functions and sets}

The Lebesgue integral turns set measure into function integration:
\[
  \int 1_E\,d\mu=\mu(E).
\]
This is conceptually important.  Sets become functions, and measuring a set becomes integrating its indicator.  Probability theory uses exactly this viewpoint: expectation is integration, and probabilities are integrals of indicators.

For the Dirichlet function,
\[
  1_{\mathbb Q\cap[a,b]}
\]
has integral zero because the rationals are countable.  For the Thomae function, which is discontinuous at every rational and continuous at every irrational, the function is still Riemann integrable because its discontinuity set is countable.  Lebesgue theory makes such examples easier to organize.

\section{Convergence theorems}

The major advantage of Lebesgue integration is not merely that it integrates more functions.  It has powerful limit theorems.

The Monotone Convergence Theorem says: if
\[
  0\le f_1\le f_2\le\cdots,
  \qquad
  f_n\to f,
\]
then
\[
  \int f\,d\mu=\lim_{n\to\infty}\int f_n\,d\mu.
\]
The Dominated Convergence Theorem says: if $f_n\to f$ pointwise, and $|f_n|\le g$ for an integrable $g$, then
\[
  \lim_{n\to\infty}\int f_n\,d\mu=\int f\,d\mu.
\]
These are theorems Riemann integration cannot match in such generality.  They explain why Lebesgue integration becomes the standard language of modern analysis.

\section{Improper Riemann integrals versus Lebesgue integrals}

A subtle point: some improper Riemann integrals converge conditionally but are not Lebesgue integrable.  For example,
\[
  \int_1^\infty \frac{\sin x}{x}\,dx
\]
converges as an improper integral, but
\[
  \int_1^\infty \left|\frac{\sin x}{x}\right|dx
\]
diverges.  Lebesgue integrability of a signed function requires integrability of the absolute value.  Thus Lebesgue integration is not simply “more permissive”; it is more structurally consistent with normed function spaces such as $L^1$.

\section{Comparison table}

\begin{center}
\small
\begin{tabularx}{\textwidth}{@{}lYY@{}}
\toprule
Aspect & Riemann & Lebesgue \\
\midrule
Partition & Domain intervals & Value levels and measurable sets \\
Basic object & Bounded functions on intervals & Measurable functions \\
Bad example & $1_{\mathbb Q}$ not integrable & Integral is $0$ \\
Limit theorems & Limited & MCT, Fatou, DCT \\
Natural space & Continuous functions and elementary calculus & $L^p$ spaces, probability, Fourier analysis \\
\bottomrule
\end{tabularx}
\end{center}


\section{Source repair: the measure-space starting point}

The uploaded note explicitly starts the Lebesgue construction from a measure space
\[
  (X,\mathcal M,\mu),
\]
where \(\mathcal M\) is a \(\sigma\)-algebra and \(\mu\) is a measure.  A function \(f:X\to\mathbb R\) is measurable if sets such as
\[
  \{x\in X:f(x)>\alpha\}
\]
belong to \(\mathcal M\) for every real \(\alpha\).  This is the point at which the theory leaves elementary calculus: before integrating, one must decide which subsets are measurable.

\section{Source repair: simple functions as building blocks}

The uploaded text defines a nonnegative simple function by
\[
  \varphi(x)=\sum_{i=1}^n a_i\chi_{A_i}(x),
\]
and sets
\[
  \int_X \varphi\,d\mu=\sum_{i=1}^n a_i\mu(A_i).
\]
For a nonnegative measurable function,
\[
  \int_X f\,d\mu
  =\sup\left\{\int_X\varphi\,d\mu:0\le\varphi\le f,\ \varphi\text{ simple}\right\}.
\]
This is the formal version of cutting by values.  Instead of approximating the domain by small intervals, we approximate the function from below by simple measurable functions.

\section{Source repair: probability viewpoint}

The source also points toward probability.  A random variable is a measurable function, and expectation is an integral:
\[
  \mathbb E[X]=\int_\Omega X\,d\mathbb P.
\]
Events are measured by probabilities, and indicators satisfy
\[
  \mathbb E[1_A]=\mathbb P(A).
\]
This makes the Lebesgue framework the natural language of modern probability, not only a generalization of calculus.

\section{A note on scope}

The slogan is useful:
\[
  \text{Riemann cuts the domain; Lebesgue cuts by values.}
\]
But the deeper difference is categorical: Riemann integration is built around approximating graphs by rectangles, while Lebesgue integration is built around measurable sets and stable limiting operations.  That is why Lebesgue integration becomes the foundation for modern analysis, probability, PDE, and functional analysis.

\section{Riemann integration partitions the domain}
Riemann integration controls oscillation on subintervals of the domain.  A bounded function is Riemann integrable when upper and lower sums can be made arbitrarily close.
\section{Lebesgue integration partitions the values}
Lebesgue integration instead measures level sets.  Simple functions approximate by values, and measurable sets replace intervals as the basic pieces.  This is why characteristic functions become fundamental.
\section{Dirichlet function as separator}
The function \(\mathbf 1_{\mathbb Q}\) is discontinuous everywhere and not Riemann integrable, but it is Lebesgue integrable with integral zero because \(\mathbb Q\) has measure zero.
\section{Convergence theorems as exchange principles}

The central question behind the convergence theorems is not only whether a limit exists.  It is whether the following exchange is legal:
\[
  \int \lim_{n\to\infty} f_n\,d\mu
  \stackrel{?}{=}
  \lim_{n\to\infty}\int f_n\,d\mu.
\]
Lebesgue theory gives precise hypotheses under which this exchange is safe.

\begin{theorem}[Monotone convergence, informal reading]
If \(0\le f_1\le f_2\le\cdots\) and \(f_n\to f\), then
\[
  \int f_n\,d\mu\longrightarrow \int f\,d\mu.
\]
The intuition is that positive mass is only accumulating upward; no cancellation or escaping negative part is hidden.
\end{theorem}

\begin{theorem}[Dominated convergence, informal reading]
If \(f_n\to f\) pointwise and there is an integrable function \(g\) such that
\[
  |f_n|\le g
\]
for all \(n\), then
\[
  \int f_n\,d\mu\longrightarrow \int f\,d\mu.
\]
The dominating function prevents mass from escaping to infinity or concentrating in narrower and narrower spikes.
\end{theorem}

Fatou's lemma is the one-sided safety rule:
\[
  \int \liminf_{n\to\infty} f_n\,d\mu
  \le
  \liminf_{n\to\infty}\int f_n\,d\mu
  \qquad (f_n\ge0).
\]
It says that, for nonnegative functions, the integral of the limiting lower mass cannot exceed the limiting lower integrals.

\begin{example}[Escaping mass]
On \(\mathbb R\), let
\[
  f_n=\mathbf 1_{[n,n+1]}.
\]
Then \(f_n(x)\to0\) for every fixed \(x\), but
\[
  \int_{\mathbb R} f_n\,dx=1
\]
for all \(n\).  Therefore
\[
  \int \lim f_n\,dx=0
  \ne
  1=\lim \int f_n\,dx.
\]
The mass has not disappeared; it has escaped to infinity.  Dominated convergence fails because no single integrable \(g\) can dominate all these translated bumps.
\end{example}

\begin{example}[Concentrating mass]
On \([0,1]\), let
\[
  f_n=n\mathbf 1_{[0,1/n]}.
\]
Then \(f_n(x)\to0\) for every \(x>0\), but
\[
  \int_0^1 f_n(x)\,dx=1.
\]
The mass concentrates into a narrower spike near \(0\).  Again, pointwise convergence alone is too weak to justify exchanging limit and integral.
\end{example}

\begin{remark}[Probability language]
In probability, the same question is whether
\[
  \mathbb E[\lim X_n]=\lim\mathbb E[X_n]
\]
is legal.  Monotone convergence, Fatou's lemma, and dominated convergence are therefore not abstract measure-theory decorations; they are rules for passing limits through expectations.
\end{remark}
\section{Integration theory through measurable structure}
Riemann integration manages oscillation on intervals; Lebesgue integration measures sets of values.  The shift from partitions of the domain to measurable level sets is the conceptual upgrade.

\section{Riemann sums and measure spaces}

The Riemann integral begins with partitions of intervals and rectangles.  Lebesgue integration changes the organizing principle: instead of cutting the domain into small intervals first, it studies the size of sets on which the function takes certain values.  A measure space is a triple
\[
  (X,\mathcal A,\mu),
\]
where $\mathcal A$ is a sigma-algebra of measurable sets and $\mu$ assigns sizes to those sets.

A simple function has the form
\[
  s=\sum_{j=1}^m c_j\mathbf 1_{E_j}.
\]
Its integral is defined by
\[
  \int s\,d\mu=\sum_{j=1}^m c_j\mu(E_j).
\]
Nonnegative measurable functions are integrated by approximating from below by simple functions, and general integrable functions are handled by positive and negative parts.

This preserves the elementary intuition of area, but it makes convergence theorems much stronger.

\section{Dominated and monotone convergence in practice}

The three convergence tools should be read as three different controls on a limiting process.

\[
\begin{array}{c|c|c}
\text{tool} & \text{hypothesis} & \text{what it prevents}\\
\hline
\text{MCT} & 0\le f_n\uparrow f & \text{loss by cancellation}\\
\text{Fatou} & f_n\ge0 & \text{overestimating the limiting lower mass}\\
\text{DCT} & |f_n|\le g\in L^1 & \text{escape or concentration of mass}
\end{array}
\]

This is why Lebesgue integration is the natural language of probability, Fourier analysis, and PDEs.  These subjects constantly form limiting objects: limits of random variables, limits of Fourier partial sums, weak or strong limits of approximate solutions.  The integral must be stable enough to survive those passages to the limit.

\section{Lp spaces and duality}

For $1\le p<\infty$, define
\[
  L^p(X)=\{f:\ \int |f|^p\,d\mu<\infty\}/\sim,
\]
where functions equal almost everywhere are identified.  The norm is
\[
  \|f\|_p=\left(\int |f|^p\,d\mu\right)^{1/p}.
\]
Hölder's inequality says that if $1/p+1/q=1$, then
\[
  \int |fg|\,d\mu\le \|f\|_p\|g\|_q.
\]
Thus integration becomes a pairing between function spaces.  The elementary integral $\int fg$ is now a duality bracket, and inequalities become statements about bounded linear functionals.

\section{Returning to the original Riemann examples}

Darboux sums, Riemann sums, and step-function pictures remain valuable because they teach integral as accumulation.  The advanced reconstruction does not erase them.  It explains why later mathematics replaces intervals by measurable sets, pointwise equality by almost-everywhere equality, and ordinary convergence by modes of convergence suitable for passing limits through integrals.
